\section{한계와 후속 연구}\label{sec:future}

\subsection{한계}
\begin{itemize}
  \item \textbf{속도.} 마지막 분할 시점 $t^*$의 상한이 없다. 정리의 속도는 $t^*$ 이후에만 적용된다. 실험의 곡선 대부분은 $t^*$ 이전이다.
  \item \textbf{가정과 구현의 차이.} 정리는 동시 갱신, RM+, 판정 구간의 균등 혼합을 가정한다. 구현은 교대 순회, 후회 매칭, bucket 평균 표를 쓴다(표~\ref{tab:gap}).
  \item \textbf{계보 비용.} 정보집합별 표 없는 순수판은 계보 비용이 분할 단계에 대해 곱으로 늘어난다. 큰 게임에는 \ref{sec:derived}절의 파생 기법이 필요하다. 그 기법들은 정리의 가정을 일부 벗어난다.
  \item \textbf{실험의 범위.} 실험은 seed 1개(일부 4–10개)이고, 정확한 exploitability는 작은 게임에서만 쟀다. 세븐 포커의 지표(LBR)는 exploitability가 아니다.
\end{itemize}

\subsection{후속 연구}
\paragraph{1. 더 나은 초기 추상화.}
\begin{itemize}
  \item \textbf{구조 서명.} 수렴에 필요한 것은 같은 플레이어, 같은 행동집합뿐이다. 그러나 앞으로의 구조까지 같은 것끼리 묶으면 필요한 분할이 줄 것으로 기대한다.
    \begin{itemize}
      \item 정보집합 $I$에서 행동 $a$ 뒤에 같은 플레이어가 다음으로 마주칠 수 있는 정보집합의 집합을 $\mathrm{Next}(I,a)$라 한다.
      \item 깊이 $d$의 구조 서명을 다음과 같이 재귀적으로 정의한다.
        \[
        \sigma_0(I)=(P(I),A(I)),\qquad\sigma_d(I)=\big(\sigma_0(I),\ (a\mapsto\{\sigma_{d-1}(J):J\in\mathrm{Next}(I,a)\})_{a}\big)
        \]
      \item 서명이 같은 것끼리 묶으면 게임 지식 없이 시뮬레이터만으로 시작 분할을 만들 수 있다.
      \item 거리를 $(1-\gamma)\sum_d\gamma^d\mathbb 1[\sigma_d(I)\ne\sigma_d(I')]$로 두면 초거리(ultrametric)가 된다. 그래서 "거리 $\le\varepsilon$"이 동치관계가 되어 그대로 bucket으로 쓸 수 있다.
    \end{itemize}
  \item \textbf{보상과 우연 확률.} 보상과 우연 확률의 차이까지 넣는 확장은 고정 추상화의 오차 상한~\citep{kroer2016imperfect,kroer2018unified}과 연결할 수 있다.
  \item \textbf{이전 실행의 재사용.} 이전 실행의 최종 분할을 다음 실행의 시작점으로 쓰는 것도 후보다.
\end{itemize}

\paragraph{2. 추상화 분할의 개선.}
\begin{itemize}
  \item \textbf{거짓 분할 없애기.} 여유 라벨 집합(보조정리~\ref{lem:slack})을 표본 판정에 넣으면 좋은 사건 위에서는 거짓 분할이 없다. 지금의 argmax와 역치 판정에는 이 성질이 없다.
  \item \textbf{숨은 후회를 직접 추정하기.} 라벨 항등식으로 $C_W(B)=\min_\ell\sum_{I\in B}e_I(\ell)$을 행동마다 선형인 합으로 바꿀 수 있다. 도달 질량에 비례해 뽑은 구성원으로 중요도 가중 추정을 하면, 표본 수가 bucket 크기와 무관해질 것으로 본다(\ref{sec:derived-theta}절). 증명은 남아 있다.
  \item \textbf{문턱의 크기.} 문턱을 bucket의 도달 질량에 비례하게 잡으면, 고정 문턱판의 오차가 최종 bucket 수에 비례하지 않게 할 수 있다.
  \item \textbf{일시적 악화 줄이기.} 1단 게임에서는 라벨이 갈린 bucket을 두 플레이어 모두 한꺼번에 나눌 때 일시적인 악화가 절반으로 줄었다. 사적 정보가 있는 게임에서 같은지 확인해야 한다.
  \item \textbf{자식 표를 어디서 시작할지.} 0에서 다시 시작할지, 부모 표를 물려줄지(regret transfer~\citep{brown2014regret}) 정해야 한다. RM+에서는 자식의 시작 노름이 부모 이하이면 보장이 유지된다.
  \item \textbf{평균 전략 방향의 치우침.} $\Omega$는 조건부 평균 표로 없앨 수 있다. $K$를 줄이는 방법은 반복 전략 몇 개를 섞는 것이다.
  \item \textbf{$t^*$의 상한.} 상위 두 행동의 누적 후회 차이에 대한 margin 조건 아래 $t^*$를 묶어 볼 수 있다.
\end{itemize}

\paragraph{3. 3인 이상의 게임.}
\begin{itemize}
  \item \textbf{성립하는 것.} 보조정리~\ref{lem:decomp}의 첫 부등식(후회 분해)과 라벨 항등식은 플레이어마다 성립한다. 그래서 ASWISS를 그대로 쓰면 각 플레이어의 평균 후회를 0으로 보내는 논증이 그대로 남을 것으로 본다.
  \item \textbf{성립하지 않는 것.} 2인 제로섬이 아니면 평균 후회가 0이라는 것이 내쉬 균형을 뜻하지 않는다. 반복 프로필들의 균등 분포가 근사 coarse correlated equilibrium이라는 데서 멈춘다.
  \item \textbf{과제.} 정리해야 할 것은 둘이다.
    \begin{itemize}
      \item 어떤 해 개념을 목표로 할지
      \item 6인 포커처럼 실전에서 쓰이는 설정~\citep{brown2019superhuman}에서 분할 규칙이 무엇을 보장하는지
    \end{itemize}
\end{itemize}

\paragraph{4. Absent-mindedness가 있을 때의 전략 개선.}
\begin{itemize}
  \item \textbf{지금 보장되는 것.} 시작 분할이나 분할 결과가 absent-minded이면(한 경로가 같은 bucket을 두 번 지나면), 원 게임 NashConv의 상한(5절)은 그대로 성립한다.
  \item \textbf{문제.} 그러나 추상 게임 안의 학습 목표가 흐려진다.
    \begin{itemize}
      \item 계획 단계의 최적과 행동 단계의 최적이 다를 수 있다~\citep{piccione1997interpretation}.
      \item bucket의 후회 합이 한 경로의 여러 방문을 따로 세기 때문이다.
    \end{itemize}
  \item \textbf{지금의 방법.} 지금은 bucket 이력 시작 칸(\ref{sec:derived-history}절)으로 absent-mindedness 자체를 피한다.
  \item \textbf{과제.} 피할 수 없는 큰 게임을 위해 두 가지가 필요하다.
    \begin{itemize}
      \item absent-minded bucket에서 일관된 행동 전략을 학습하는 갱신 규칙
      \item 그 규칙이 원 게임 상한과 어떻게 맞물리는지에 대한 분석
    \end{itemize}
\end{itemize}

\paragraph{5. 관측 분할(옛 prior 추상화)의 구성.}
\begin{itemize}
  \item \textbf{이름.} 설계 문서의 "prior 추상화"는 전략을 보기 전에 관측만으로 정하는 분할이라 붙은 이름이다. 베이즈 사전분포와 혼동되므로 이 논문은 \emph{관측 분할}이라 부른다(\ref{sec:derived-pp}절).
  \item \textbf{지금 방식의 약점.} 지금의 관측 분할은 무작위 후보 자르기(특징, 문턱값)를 시험해 나눈다. 후보가 무작위라서 필요한 자르기를 찾기까지 오래 걸린다.
  \item \textbf{개선 후보.}
    \begin{itemize}
      \item 분할 전 전략으로 얻은 표본의 후회 벡터나 가치 벡터 사이의 거리로 라벨을 붙인다. 그 라벨을 맞히는 결정 트리를 학습해 자르기를 고른다.
      \item 학습한 관측 임베딩이나 카드 codebook을 특징으로 쓴다.
      \item 관측 분할이 세분된 뒤 포인터를 다시 배정하는 탐색 기간을 둔다.
      \item 잎 수에 상한을 둔다.
    \end{itemize}
  \item \textbf{필요한 분석.} 관측 분할은 표를 바꾸지 않으므로 수렴 보장과 독립적으로 고칠 수 있다. 다만 관측 분할이 충분히 세밀하지 않으면 전략 분할이 나눌 수 없는 숨은 후회가 남는다. 그 크기에 대한 분석이 필요하다.
\end{itemize}

\paragraph{6. 그 밖의 과제.}
\begin{itemize}
  \item 교대 갱신과 선형 평균으로 증명을 넓힌다.
  \item 저장한 반복 전략의 혼합으로 정리의 대상인 평균 정책을 내놓는다.
  \item seed 반복과 신뢰구간을 낸다.
  \item 같은 크기와 같은 정보를 가진 고정 추상화와 비교한다.
  \item 신경망 기준선을 GPU에서 비교한다.
\end{itemize}
